\chapter{Bridges from finite sequences to spectra}\label{ch:bridges-signals}
\begin{readergoal}Establish why discrete Fourier coefficients, filtering, power, phase, and aliasing behave as they do before treating recorded LFP samples as continuous-time signals.\end{readergoal}

\bridgeentry{fourier-orthogonality}{Complex exponentials form an orthogonal finite family}
\begin{proposition}For integers $k,\ell\in\{0,\ldots,N-1\}$, $N\geq1$,
\[\sum_{n=0}^{N-1}e^{2\pi i(k-\ell)n/N}=\begin{cases}N,&k=\ell,\\0,&k\ne\ell.\end{cases}\]\end{proposition}
\begin{proof}For $k=\ell$, every term is one. Otherwise set $r=e^{2\pi i(k-\ell)/N}$; then $r^N=1$ but $r\ne1$. The finite geometric sum equals $(1-r^N)/(1-r)=0$.\end{proof}
\paragraph{Example.} This finite orthogonality is the algebraic reason a DFT separates integer-grid Fourier bins; an off-grid sinusoid leaks across bins under a finite window.

\bridgeentry{dft-inverse}{The DFT admits an exact finite inverse}
\begin{proposition}For $x_0,\ldots,x_{N-1}\in\mathbb C$, set $X_k=\sum_{n=0}^{N-1}x_ne^{-2\pi i kn/N}$. Then
\[x_n=\frac1N\sum_{k=0}^{N-1}X_ke^{2\pi i kn/N}.\]\end{proposition}
\begin{proof}Insert the definition of $X_k$ into the inverse expression and exchange two \emph{finite} sums. The coefficient of $x_m$ is $N^{-1}\sum_k e^{2\pi i k(n-m)/N}$, which equals one for $m=n$ and zero otherwise by the previous bridge. Only $x_n$ remains.\end{proof}
\paragraph{Example.} An inverse DFT exactly recovers the original finite samples up to numerical roundoff. It does not prove recovery of the unknown continuous biological signal between sample times.

\bridgeentry{parseval-dft}{Finite Fourier energy has an exact accounting identity}
\begin{proposition}With the DFT normalization above, $\sum_n|x_n|^2=\frac1N\sum_k|X_k|^2$.\end{proposition}
\begin{proof}Insert the inverse DFT into $\sum_n x_n\bar x_n$. Expand the finite sums and group terms by pairs $(k,\ell)$. The inner sum over $n$ vanishes whenever $k\ne\ell$ by orthogonality and equals $N$ when $k=\ell$, leaving $N^{-1}\sum_k|X_k|^2$.\end{proof}
\paragraph{Example.} Periodogram normalization should respect energy units and sampling interval. ``Power'' and ``power spectral density'' differ by frequency-bin widths.

\bridgeentry{circular-convolution}{Circular convolution becomes multiplication under the DFT}
\begin{proposition}If $(x\circledast h)_n=\sum_{m=0}^{N-1}x_mh_{(n-m)\bmod N}$, then $\operatorname{DFT}(x\circledast h)_k=X_kH_k$.\end{proposition}
\begin{proof}Substitute the circular sum into the DFT, exchange finite sums, and for each $m$ set $j=(n-m)\bmod N$. Periodicity gives $e^{-2\pi ikn/N}=e^{-2\pi ikm/N}e^{-2\pi ikj/N}$. The two sums factor into $X_k$ and $H_k$.\end{proof}
\paragraph{Example.} Multiplying two equal-length FFTs implements \emph{circular} convolution. To obtain linear finite convolution one must pad sufficiently or handle boundaries separately; otherwise opposite edges mix.

\bridgeentry{shift-phase}{Time shifts rotate Fourier phase}
\begin{proposition}For a circular shift $y_n=x_{(n-n_0)\bmod N}$, the DFT is $Y_k=e^{-2\pi ikn_0/N}X_k$.\end{proposition}
\begin{proof}In the DFT sum set $m=(n-n_0)\bmod N$. Periodicity preserves the summation domain and extracts the multiplier $e^{-2\pi ikn_0/N}$. The multiplier has modulus one, so magnitudes are preserved and phases change.\end{proof}
\paragraph{Example.} Two copies of the same oscillation with a constant temporal shift can have identical PSD yet different cross-spectral phase. PSD alone cannot recover relative timing.

\bridgeentry{aliasing-equivalence}{Uniform samples cannot distinguish aliased sinusoids}
\begin{proposition}For sample rate $F_s>0$, every $f\in\mathbb R$ and integer $m$ satisfy
\[e^{2\pi i (f+mF_s)n/F_s}=e^{2\pi i fn/F_s}\quad\text{for all integers }n.\]\end{proposition}
\begin{proof}The ratio of the two exponentials equals $e^{2\pi i mn}=1$ for every integer $n$. Hence their sample sequences coincide exactly, although their continuous frequencies differ.\end{proof}
\paragraph{Example.} A 90-Hz tone sampled at 100 Hz can appear at a frequency of 10 Hz after folding; no digital filter applied afterward can establish which continuous sinusoid generated the samples. Acquisition requires appropriate analog anti-aliasing.

\begin{exercises}
\item Work out the DFT and inverse DFT of $[1,0,0,0]$ by hand.
\item Numerically compare circular convolution with linear convolution for a short impulse near an array boundary.
\item Demonstrate DFT magnitude invariance under circular shifts and the corresponding phase change.
\end{exercises}